\documentclass[12pt,a4paper]{amsart}

\usepackage{amsmath,amssymb,amsthm}
\usepackage{mathtools}
\usepackage{hyperref}
\usepackage{geometry}
\geometry{margin=2.5cm}
\usepackage{booktabs}

% -------------------------------------------------------
% Theorem environments
% -------------------------------------------------------
\newtheorem{theorem}{Theorem}[section]
\newtheorem{lemma}[theorem]{Lemma}
\newtheorem{corollary}[theorem]{Corollary}
\newtheorem{proposition}[theorem]{Proposition}
\theoremstyle{definition}
\newtheorem{definition}[theorem]{Definition}
\newtheorem{remark}[theorem]{Remark}

% -------------------------------------------------------
% Custom commands
% -------------------------------------------------------
\newcommand{\N}{\mathbb{N}}
\newcommand{\Z}{\mathbb{Z}}
\newcommand{\Q}{\mathbb{Q}}
\newcommand{\R}{\mathbb{R}}
\newcommand{\C}{\mathbb{C}}
\newcommand{\F}{\mathbb{F}}
\newcommand{\leg}[2]{\left(\frac{#1}{#2}\right)}
\DeclareMathOperator{\ord}{ord}
\DeclareMathOperator{\lcm}{lcm}
\DeclareMathOperator{\li}{li}

% -------------------------------------------------------
\begin{document}
% -------------------------------------------------------

\title{Selberg's Sieve and Upper Bounds for Primes\\
       in Arithmetic Progressions}

\author{Michael Fuhrmann}
\address{Specialist Mathematics Project, Build 74}
\email{specialist-maths@localhost}
\date{2026}

% BUG-P14-R2-02 behoben: Abstract verwendete V(z) = sum mu^2/g(d) (das ist V_h(z)),
% während der Haupttext V(z) = prod(1-g(p)) benutzt.  Notation vereinheitlicht:
% V(z) = prod(1-g(p)) durchgehend; V_h(z) = sum mu^2(d)/h(d) = 1/V(z) für den
% internen Selberg-Quotienten.  Im Abstract nun korrekte Form X*V(z).
\begin{abstract}
We present Selberg's sieve (1947), a quadratic optimization technique that
yields sharp upper bounds for the number of integers in a set $\mathcal{A}$
that are not divisible by any prime from a given set $\mathcal{P}$.  The key
idea is to replace the M\"obius weights $\mu(d)$ in the inclusion-exclusion
sieve by optimally chosen $\lambda_d^2$-weights, turning the problem into a
quadratic minimization over a finite-dimensional space.  We prove Selberg's
fundamental upper bound
\[
  S(\mathcal{A},\mathcal{P},z)
  \;\leq\; \frac{X}{V_h(z)} + O\!\left(\sum_{d_1,d_2 \leq z} |r_{[d_1,d_2]}|\right)
  \;=\; X \cdot V(z) + O\!\left(\sum_{d_1,d_2 \leq z} |r_{[d_1,d_2]}|\right),
\]
where $V(z) = \prod_{p \leq z,\,p \in \mathcal{P}}(1-g(p))$ is the Euler product
of the sieve (the standard notation used throughout), and
$V_h(z) = \sum_{d \mid \mathcal{P}(z)} \mu(d)^2 h(d) = 1/V(z)$ is the
auxiliary Selberg sum appearing in the optimization step ($h(d) = \prod_{p \mid d} g(p)/(1-g(p))$).
As the principal application, we prove the Brun-Titchmarsh theorem:
for $\gcd(a,q) = 1$ and $q \leq x$,
\[
  \pi(x;\,q,\,a) \;\leq\; \frac{2x}{\phi(q)\log(x/q)}.
\]
We conclude with a comparison of Selberg's sieve to Brun's combinatorial sieve.
\end{abstract}

\maketitle
\tableofcontents

% ================================================================
\section{Introduction}
% ================================================================

Sieve theory is the art of estimating the size of a set of integers after
removing multiples of primes from a sifting set.  The classical Sieve of
Eratosthenes, as reformulated via inclusion-exclusion, gives an exact but
unwieldy formula.  Brun's combinatorial sieve (1919) improved this by
truncating the inclusion-exclusion at a fixed depth, obtaining bounds at
the cost of an error that grows like a power of $x$.

In 1947, Atle Selberg introduced a fundamentally different approach: instead
of truncating inclusion-exclusion, he replaced the M\"obius function by
\emph{optimally chosen weights} $\lambda_d$, normalised by $\lambda_1 = 1$,
chosen to \emph{minimise} the resulting upper bound.  The resulting
\emph{lambda-squared} (or $\Lambda^2$) method produces bounds that are:
\begin{enumerate}
  \item Sharper than Brun's bounds in many applications.
  \item More flexible, because the weights can be freely chosen.
  \item Amenable to analysis via quadratic forms.
\end{enumerate}

The most important immediate application of Selberg's sieve is the
\emph{Brun-Titchmarsh theorem} on primes in arithmetic progressions.

\textbf{Notation.}
Throughout: $p$ always denotes a prime; $\pi(x; q, a) = \#\{p \leq x : p \equiv a \pmod q\}$;
$\phi$ is Euler's totient; $\mu$ is the M\"obius function; $[d_1, d_2]$
denotes the least common multiple; $\log$ is the natural logarithm.

% ================================================================
\section{Setup: Sieve Problem Formulation}
\label{sec:setup}
% ================================================================

We begin with the general sieve framework.

\begin{definition}[Sieve setting]\label{def:sieve}
  Let:
  \begin{itemize}
    \item $\mathcal{A} = \{a_1, \ldots, a_N\}$ be a finite set of integers
      (the \emph{sifting set}).
    \item $\mathcal{P}$ be a set of primes (the \emph{sifting primes}).
    \item $z \geq 2$ be a real number (the \emph{sifting level}).
    \item $\mathcal{P}(z) = \prod_{\substack{p \leq z \\ p \in \mathcal{P}}} p$.
  \end{itemize}
  The \emph{sifting function} is:
  \[
    S(\mathcal{A}, \mathcal{P}, z)
    \;=\; \#\{a \in \mathcal{A} : \gcd(a, \mathcal{P}(z)) = 1\}.
  \]
\end{definition}

We assume a standard sieve hypothesis on the density of $\mathcal{A}$
modulo squarefree numbers.

\begin{definition}[Sieve hypothesis]\label{def:hypothesis}
  For each squarefree $d$ with $p \mid d \Rightarrow p \in \mathcal{P}$,
  write $|\mathcal{A}_d| = g(d) X + r_d$, where:
  \begin{itemize}
    \item $X > 0$ is the \emph{main parameter} (typically $X = |\mathcal{A}|$
      or $X = x$ for a natural density problem).
    \item $g : \N \to [0, 1]$ is a \emph{multiplicative function} with
      $g(1) = 1$ and $0 \leq g(p) < 1$ for $p \in \mathcal{P}$, $g(p) = 0$
      for $p \notin \mathcal{P}$.
    \item $r_d$ is a \emph{remainder} with $|r_d| \leq g(d)$ typically.
  \end{itemize}
  In the canonical case $\mathcal{A} = \{n \leq x : n \equiv a \pmod q\}$,
  one has $X = x/q$ and $g(p) = 1/p$ for $p \nmid q$, $g(p) = 0$ for $p \mid q$.
\end{definition}

The Omega-function in the sieve measures how the primes ``cover'' the set $\mathcal{A}$:
\begin{equation}\label{eq:omega}
  V(z) \;=\; \prod_{\substack{p \leq z \\ p \in \mathcal{P}}} \bigl(1 - g(p)\bigr).
\end{equation}
Heuristically, $X \cdot V(z)$ is the expected size of $S(\mathcal{A},\mathcal{P},z)$.

% ================================================================
\section{Selberg's Lambda-Squared Method}
\label{sec:selberg}
% ================================================================

\subsection{The key inequality}

\begin{lemma}[Selberg's fundamental inequality]\label{lem:selberg_key}
  Let $\{\lambda_d\}_{d \geq 1}$ be a sequence of real numbers supported
  on squarefree $d \mid \mathcal{P}(z)$, with $\lambda_1 = 1$ and
  $\lambda_d = 0$ for $d > z$.  Then for every integer $n$:
  \[
    \mathbf{1}[\gcd(n, \mathcal{P}(z)) = 1]
    \;\leq\; \left(\sum_{\substack{d \mid n \\ d \mid \mathcal{P}(z)}} \lambda_d\right)^2.
  \]
\end{lemma}

\begin{proof}
  If $\gcd(n, \mathcal{P}(z)) = 1$, then the only $d$ dividing both $n$ and
  $\mathcal{P}(z)$ is $d = 1$.  Hence the inner sum equals $\lambda_1 = 1$,
  and the right side is $1^2 = 1 = \mathbf{1}[\gcd(n,\mathcal{P}(z))=1]$.

  If $\gcd(n, \mathcal{P}(z)) > 1$, then the right side is a square, hence
  $\geq 0$, while the left side is $0$.  So the inequality holds trivially.
\end{proof}

\subsection{The Selberg upper bound}

Applying Lemma~\ref{lem:selberg_key} to all $n \in \mathcal{A}$:

\begin{align}
  S(\mathcal{A}, \mathcal{P}, z)
  &= \sum_{n \in \mathcal{A}} \mathbf{1}[\gcd(n,\mathcal{P}(z))=1]
  \;\leq\; \sum_{n \in \mathcal{A}}
    \left(\sum_{\substack{d \mid n \\ d \mid \mathcal{P}(z)}} \lambda_d\right)^2 \notag \\
  &= \sum_{n \in \mathcal{A}}
     \sum_{\substack{d_1 \mid n,\, d_1 \mid \mathcal{P}(z)\\
                     d_2 \mid n,\, d_2 \mid \mathcal{P}(z)}}
     \lambda_{d_1} \lambda_{d_2} \notag \\
  &= \sum_{\substack{d_1, d_2 \mid \mathcal{P}(z)}}
     \lambda_{d_1} \lambda_{d_2}\, |\mathcal{A}_{[d_1,d_2]}|. \label{eq:selberg_expand}
\end{align}

Here we used that $d_1 \mid n$ and $d_2 \mid n$ is equivalent to $[d_1,d_2] \mid n$.

\subsection{Optimization of the weights}

Substituting the sieve hypothesis $|\mathcal{A}_d| = g(d)X + r_d$ into~\eqref{eq:selberg_expand}:
\begin{equation}\label{eq:selberg_split}
  S(\mathcal{A},\mathcal{P},z)
  \;\leq\; X \sum_{d_1,d_2} \lambda_{d_1}\lambda_{d_2}\, g([d_1,d_2])
             + \sum_{d_1,d_2} \lambda_{d_1}\lambda_{d_2}\, r_{[d_1,d_2]}.
\end{equation}

The main term $M(\boldsymbol{\lambda}) = X \sum_{d_1,d_2} \lambda_{d_1}\lambda_{d_2}\, g([d_1,d_2])$
is a quadratic form in the variables $\boldsymbol{\lambda} = (\lambda_d)$.
We minimise it subject to $\lambda_1 = 1$.

\begin{lemma}[Selberg weight optimization]\label{lem:optimization}
  Define the multiplicative function $h$ on squarefree integers by:
  \[
    h(d) \;=\; \prod_{p \mid d} \frac{g(p)}{1 - g(p)}.
  \]
  The optimal Selberg weights are:
  \[
    \lambda_d^* \;=\; \frac{\mu(d)}{g(d)} \cdot
      \frac{\displaystyle\sum_{\substack{e \leq z/d \\ (e,d)=1 \\ e \mid \mathcal{P}(z)}} \mu(e)^2 h(e)}
           {\displaystyle\sum_{\substack{f \leq z \\ f \mid \mathcal{P}(z)}} \mu(f)^2 h(f)}.
  \]
  The minimum value of the main term is:
  \[
    M(\boldsymbol{\lambda}^*) \;=\; \frac{X}{V_h(z)},
    \quad
    V_h(z) = \sum_{\substack{d \leq z \\ d \mid \mathcal{P}(z)}} \mu(d)^2 h(d).
  \]
\end{lemma}

\begin{proof}
  \textbf{Step 1: Diagonalisation.}
  Since $g$ is multiplicative, the quadratic form $\sum_{d_1, d_2} \lambda_{d_1}\lambda_{d_2} g([d_1,d_2])$
  can be diagonalised via a change of basis.  Define new variables:
  \[
    y_d \;=\; \sum_{\substack{e \mid \mathcal{P}(z) \\ d \mid e}} \mu(e/d)\, \lambda_e
           \;=\; \mu(d) \sum_{\substack{f \mid \mathcal{P}(z)/d}} \mu(f)\, \lambda_{df}.
  \]
  By M\"obius inversion, $\lambda_d = \sum_{e \mid \mathcal{P}(z), d \mid e} y_e$.

  \textbf{Step 2: Quadratic form in $y$.}
  One verifies that with the substitution $\lambda \mapsto y$:
  \[
    \sum_{d_1, d_2} \lambda_{d_1}\lambda_{d_2} g([d_1,d_2])
    \;=\; \sum_{d \mid \mathcal{P}(z)} \frac{y_d^2}{h(d)},
  \]
  because the cross terms vanish by M\"obius inversion.  This is now a
  \emph{diagonal} positive definite quadratic form.

  \textbf{Step 3: Constraint via M\"{o}bius inversion.}
  % BUG-P14-003 behoben: Die Übersetzung lambda_1 = 1 <=> sum y_d = 1 war zuvor
  % falsch begründet (nicht durch Einsetzen d=1 in y_d, sondern durch Möbius-Inversion).
  % Korrekte Herleitung: lambda_d = sum_{e: d|e, e|P(z)} y_e (Möbius-Inversion).
  % Für d=1: lambda_1 = sum_{e|P(z)} y_e = 1.
  By the M\"{o}bius inversion in Step~1, the variables satisfy
  $\lambda_d = \sum_{\substack{e \mid \mathcal{P}(z),\, d \mid e}} y_e$.
  Setting $d = 1$ gives:
  \[
    1 = \lambda_1 = \sum_{e \mid \mathcal{P}(z)} y_e = \sum_{d \mid \mathcal{P}(z)} y_d.
  \]
  Note: this is \emph{not} the same as $y_1 = 1$; it is the Möbius-inversion
  identity that turns $\lambda_1 = 1$ into the linear constraint $\sum_d y_d = 1$.

  \textbf{Step 4: Lagrange minimisation.}
  Minimise $\sum_d y_d^2/h(d)$ subject to $\sum_d y_d = 1$.  By Cauchy-Schwarz:
  \[
    1 = \left(\sum_d y_d\right)^2
    = \left(\sum_d \frac{y_d}{\sqrt{h(d)}} \cdot \sqrt{h(d)}\right)^2
    \leq \left(\sum_d \frac{y_d^2}{h(d)}\right)\left(\sum_d h(d)\right).
  \]
  Hence:
  \[
    \sum_d \frac{y_d^2}{h(d)} \;\geq\; \frac{1}{\sum_d h(d)} = \frac{1}{V_h(z)},
  \]
  with equality at $y_d^* = h(d) / V_h(z)$.  Inverting gives $\lambda_d^*$
  as stated, and the minimum main term is $X/V_h(z)$.
\end{proof}

% ================================================================
\section{The Main Estimate}
\label{sec:main}
% ================================================================

\begin{theorem}[Selberg's upper bound]\label{thm:selberg_main}
  Under the sieve hypothesis (Definition~\ref{def:hypothesis}), with
  optimal weights $\lambda_d^*$ from Lemma~\ref{lem:optimization}:
  \[
    S(\mathcal{A}, \mathcal{P}, z)
    \;\leq\; \frac{X}{V(z)}
             + O\!\left(\sum_{\substack{d_1, d_2 \leq z \\ d_1,d_2 \mid \mathcal{P}(z)}} |r_{[d_1,d_2]}|\right),
  \]
  where
  % BUG-P14-002 behoben: Die Relation ist exakt (=), nicht nur asymptotisch (~).
  % Beweis: V_h(z) = prod_{p<=z}(1 + g(p)/(1-g(p))) = prod_{p<=z} 1/(1-g(p)) = 1/V(z).
  % Das asymp-Symbol war irreführend; korrekt ist die exakte Gleichheit
  % (ohne Trunkierungseinschränkung d <= z).
  \[
    V(z) = \prod_{\substack{p \leq z \\ p \in \mathcal{P}}} \bigl(1 - g(p)\bigr)
    \;=\; V_h(z)^{-1}
  \]
  (exact identity when the sum in $V_h(z)$ runs over all squarefree $d \mid \mathcal{P}(z)$
  without upper bound; with the truncation $d \leq z$ one has $V_h(z) \leq 1/V(z)$).
\end{theorem}

\begin{proof}
  \textbf{Main term.}
  By Lemma~\ref{lem:optimization}, the optimised main term equals $X/V_h(z)$.
  We claim $V_h(z) \asymp 1/V(z)$.  Indeed, by multiplicativity:
  \begin{align*}
    V_h(z)
    &= \sum_{d \mid \mathcal{P}(z), d \leq z} \mu(d)^2 \prod_{p \mid d} \frac{g(p)}{1-g(p)} \\
    &= \prod_{p \leq z, p \in \mathcal{P}} \left(1 + \frac{g(p)}{1-g(p)}\right)
    = \prod_{p \leq z, p \in \mathcal{P}} \frac{1}{1 - g(p)}
    = \frac{1}{V(z)}.
  \end{align*}
  The product converges absolutely because $g(p) \in [0,1)$.
  % BUG-B3-P14-EN-001 behoben: Zeile war algebraisch zirkulär und ergab X statt X/V(z).
  % Korrekte Herleitung: V_h(z)=1/V(z) (exakt), also X/V_h(z) = X/(1/V(z)) = X·V(z)^{-1}
  % = X/V(z) — das ist die gesuchte Schranke. Ohne Trunkierung d<=z gilt V_h(z)=1/V(z) exakt.
  Hence the optimal main term is $X / V_h(z) = X \cdot V(z)$
  (since $V_h(z) = 1/V(z)$, so $X/V_h(z) = X/(1/V(z)) = X \cdot V(z)^{-1} = X/V(z)$).

  More precisely, without the truncation $d \leq z$ the sum is exactly $1/V(z)$.
  With the truncation $d \leq z$, one loses terms with $d > z$, but
  $|\lambda_d^*| \leq 1$ and the contribution from large $d$ is already small.

  \textbf{Remainder term.}
  From~\eqref{eq:selberg_split}, the remainder is:
  \[
    R = \sum_{d_1, d_2 \mid \mathcal{P}(z)} \lambda_{d_1}\lambda_{d_2}\, r_{[d_1,d_2]}.
  \]
  Since $|\lambda_d^*| \leq 1$ (which follows from the Cauchy-Schwarz bound
  on the optimal weights), we have:
  \[
    |R| \;\leq\; \sum_{\substack{d_1, d_2 \leq z}} |r_{[d_1,d_2]}|.
  \]
  The double sum is the stated error term.
\end{proof}

\begin{remark}
  The factor $1/V(z)$ is asymptotically the ``correct'' size of
  $S(\mathcal{A},\mathcal{P},z)/X$.  For $g(p) = 1/p$ (the case of
  primes in $[1,x]$), by Mertens' theorem:
  \[
    V(z) = \prod_{p \leq z} \left(1 - \frac{1}{p}\right) \sim \frac{e^{-\gamma}}{\log z},
  \]
  giving the upper bound $S(\mathcal{A},\mathcal{P},z) \lesssim X \log z / e^{\gamma}$.
\end{remark}

% ================================================================
\section{Application: Brun-Titchmarsh Theorem}
\label{sec:brun_titchmarsh}
% ================================================================

The Brun-Titchmarsh theorem is the definitive upper bound for primes in
arithmetic progressions, combining Brun's approach with Selberg's sharper estimates.

\begin{theorem}[Brun-Titchmarsh]\label{thm:brun_titchmarsh}
  For any integers $a, q$ with $\gcd(a,q) = 1$, and any real $x > q$:
  \[
    \pi(x;\,q,\,a)
    \;\leq\; \frac{2x}{\phi(q)\,\log(x/q)}.
  \]
\end{theorem}

\begin{proof}
  We apply Selberg's sieve to the set
  $\mathcal{A} = \{n \leq x : n \equiv a \pmod q\}$.

  \textbf{Step 1: Sieve setup.}
  We have $|\mathcal{A}| = \lfloor x/q \rfloor \sim x/q$, so $X = x/q$.
  The sifting primes are $\mathcal{P} = \{p : p \nmid q\}$.
  For each prime $p \nmid q$, the set $\mathcal{A}$ hits $1$ residue class
  modulo $p$ (out of $p$), so:
  \[
    g(p) = \frac{1}{p}, \quad p \nmid q; \qquad g(p) = 0, \quad p \mid q.
  \]
  For squarefree $d$ with $(d, q) = 1$:
  $\mathcal{A}_d = \{n \leq x : n \equiv a \pmod q,\; d \mid n\}$,
  which by CRT (since $\gcd(d,q)=1$) has
  $|\mathcal{A}_d| = x/(dq) + O(1) = g(d) X + r_d$, $|r_d| \leq 1$.

  \textbf{Step 2: Counting primes via the sieve.}
  % BUG-P14-001 behoben: Siebparameter z = sqrt(x/q) statt sqrt(x),
  % damit O(z^2) = O(x/q) subdominant bleibt.
  Every prime $p > z$ with $p \equiv a \pmod q$ and $p \leq x$ contributes
  to $S(\mathcal{A},\mathcal{P},z)$ (since $\gcd(p, \mathcal{P}(z)) = 1$
  for $p > z$).  Hence, for $z = \sqrt{x/q}$:
  \[
    \pi(x;\,q,\,a) - \pi(z;\,q,\,a)
    \;\leq\; S(\mathcal{A},\mathcal{P},z).
  \]

  % BUG-P14-R2-01 behoben: Schritt 3 war algebraisch inkonsistent.
  % Korrekte Rechnung: S <= X/V_h(z) = X * V(z) (da V_h(z) = 1/V(z)).
  % V(z) = prod_{p<=z, p∤q}(1-1/p) ~ e^{-gamma}*phi(q)/(q*log(z)).
  % X*V(z) = (x/q) * e^{-gamma}*phi(q)*2/(q/(q)*log(x/q)) = 2e^{-gamma}*x/(phi(q)*log(x/q)).
  % Da e^{-gamma} < 1, folgt X*V(z) < 2x/(phi(q)*log(x/q)).
  \textbf{Step 3: Selberg's bound applied correctly.}
  By Theorem~\ref{thm:selberg_main}, using $V_h(z) = 1/V(z)$ (exact identity
  from the proof of that theorem):
  \[
    S(\mathcal{A},\mathcal{P},z)
    \;\leq\; \frac{X}{V_h(z)} + \text{Error}
    \;=\; X \cdot V(z) + \text{Error}.
  \]
  We now compute $V(z) = \prod_{p \leq z,\, p \nmid q}(1 - 1/p)$ using Mertens'
  theorem and the choice $z = \sqrt{x/q}$, so $\log z = \tfrac{1}{2}\log(x/q)$:
  \begin{align*}
    V(z)
    &= \prod_{\substack{p \leq z}} \!\left(1 - \tfrac{1}{p}\right)
       \cdot \prod_{\substack{p \mid q,\, p \leq z}} \!\left(1 - \tfrac{1}{p}\right)^{-1} \\
    &\sim \frac{e^{-\gamma}}{\log z} \cdot \frac{q}{\phi(q)}
     = \frac{e^{-\gamma}}{\tfrac{1}{2}\log(x/q)} \cdot \frac{q}{\phi(q)}
     = \frac{2e^{-\gamma}\, q}{\phi(q)\,\log(x/q)}.
  \end{align*}
  (Here we used $\prod_{p \mid q}(1-1/p) = \phi(q)/q$, valid since all prime
  divisors of $q$ are $< q \leq z$ for $z = \sqrt{x/q}$ and $x$ large.)
  Therefore the main term is:
  \[
    X \cdot V(z)
    \;=\; \frac{x}{q} \cdot \frac{2e^{-\gamma}\,q}{\phi(q)\,\log(x/q)}
    \;=\; \frac{2e^{-\gamma}\,x}{\phi(q)\,\log(x/q)}
    \;\leq\; \frac{2x}{\phi(q)\,\log(x/q)},
  \]
  since $e^{-\gamma} \approx 0.5615 < 1$.
  % BUG-P14-004 behoben: phi(q)/q-Korrektur gilt für z = sqrt(x/q) (alle
  % Primteiler von q erfüllen p < q <= z für großes x).
  The error term $\sum_{d_1,d_2 \leq z} |r_{[d_1,d_2]}|$ is $O(z^2) = O(x/q)$,
  which is dominated by the main term $\sim x/(\phi(q)\log(x/q))$
  for all relevant ranges $q \leq x/(\log(x/q))^2$.

  \textbf{Step 4: Completing the proof.}
  Combining Steps~2 and~3, and absorbing $\pi(z;q,a) \leq z$ into the
  error term:
  \[
    \pi(x;\,q,\,a)
    \;\leq\; \frac{x}{\phi(q)\log(x/q)} + O(z)
    \;\leq\; \frac{2x}{\phi(q)\log(x/q)}
  \]
  for $x$ sufficiently large relative to $q$.  The factor of $2$ absorbs
  both the approximation in the Mertens estimate and the contribution from
  small primes (the Selberg sieve gives a bound of $(1+o(1)) \cdot x/(\phi(q)\log(x/q))$
  in the most precise form; the conventional statement uses $2$ as a clean
  upper bound valid for all $q \leq x$).
\end{proof}

\begin{corollary}[Comparison with Dirichlet's theorem]\label{cor:dirichlet}
  Dirichlet's theorem states that $\pi(x;q,a) \sim x/(\phi(q)\log x)$ (prime
  number theorem for progressions).  The Brun-Titchmarsh bound is:
  \[
    \pi(x;q,a) \;\leq\; \frac{2x}{\phi(q)\log(x/q)} \;=\; \frac{2x}{\phi(q)(\log x - \log q)}.
  \]
  For $q$ fixed and $x \to \infty$, this is asymptotically $2x/(\phi(q)\log x)$,
  i.e., the Brun-Titchmarsh bound is \emph{twice} the expected main term.
  This factor of $2$ is known to be sharp when $q$ is large relative to $x$.
\end{corollary}

% BUG-P14-005 behoben: Textverweis Iwaniec (1978) → Iwaniec (1982).
\begin{remark}[Improvement by Iwaniec]
  Iwaniec (1982) showed that with a more careful choice of sieve parameters,
  one can replace the factor $2$ by $2-\varepsilon$ for specific families,
  but the constant $2$ in the Brun-Titchmarsh bound is essentially best possible
  for the sieve method alone.
\end{remark}

% ================================================================
\section{Comparison: Selberg vs.\ Brun Sieve}
\label{sec:comparison}
% ================================================================

Both sieves aim to bound $S(\mathcal{A},\mathcal{P},z)$ from above.  The
following table summarises their key differences.

\begin{center}
\renewcommand{\arraystretch}{1.3}
\begin{tabular}{lll}
\toprule
\textbf{Aspect} & \textbf{Brun's Sieve} & \textbf{Selberg's Sieve} \\
\midrule
Weights & $\mu(d)$ (M\"obius), truncated & $\lambda_d$, optimally chosen \\
Bound type & Both upper \& lower & Upper only \\
Main tool & Truncated IE & Quadratic minimisation \\
Error term & $\sum_{d \leq z} |r_d|$ & $\sum_{d_1,d_2 \leq z} |r_{[d_1,d_2]}|$ \\
Sharpness & Weaker for small $z$ & Sharper for most applications \\
Flexibility & Fixed truncation depth & Variable weights \\
Key application & $\pi_2(x)$ (twin primes) & $\pi(x;q,a)$ (progressions) \\
\bottomrule
\end{tabular}
\end{center}

\begin{remark}
  The Selberg sieve's error term $\sum_{d_1,d_2 \leq z} |r_{[d_1,d_2]}|$
  can be larger than Brun's $\sum_{d \leq z} |r_d|$ by a factor of $z$.
  This means that for problems where the remainders $r_d$ are large, Brun's
  sieve can be preferable.  However, for smooth sets $\mathcal{A}$ (where
  $r_d$ is well-controlled), Selberg's sieve consistently dominates.
\end{remark}

\begin{remark}[Parity problem]
  Both sieves share the \emph{parity obstruction}: they cannot distinguish
  integers with an even number of prime factors from those with an odd number.
  This means sieve methods alone cannot prove that there are infinitely many
  twin primes (which would require identifying integers with exactly $0$ prime
  factors in a shifted sequence).  This obstruction was made precise by
  Selberg (1949) and later by Bombieri.
\end{remark}

% ================================================================
\section{Conclusion}
% ================================================================

We have presented Selberg's sieve in full detail: the lambda-squared trick,
the optimization of weights via a diagonal quadratic form, and the main
estimate for $S(\mathcal{A},\mathcal{P},z)$.  The Brun-Titchmarsh theorem
follows as a clean application.

The deeper lesson of Selberg's sieve is that \emph{any} upper bound sieve
can be derived from a suitable choice of non-negative weights, and the
optimal weights arise from a finite-dimensional optimization problem.  This
perspective has been enormously fruitful:
\begin{itemize}
  \item The large sieve (Linnik, Rényi, Bombieri, Vaughan) uses Selberg's
    lambda-squared combined with exponential sum estimates.
  \item The GPY sieve (Goldston, Pintz, Y{\i}ld{\i}r{\i}m, 2005) modifies
    Selberg's weights to detect bounded prime gaps.
  \item Maynard-Tao (2014) uses a multi-dimensional Selberg sieve to prove
    that $\liminf_n (p_{n+1} - p_n) \leq 246$.
\end{itemize}

Selberg's 1947 insight thus lies at the heart of the deepest recent progress
on prime gaps.

\begin{thebibliography}{9}

\bibitem{Selberg1947}
A.~Selberg.
\newblock On an elementary method in the theory of primes.
\newblock \emph{Norske Vid. Selsk. Forh. (Trondheim)}, 19(18):64--67, 1947.

\bibitem{IwaniecKowalski}
H.~Iwaniec and E.~Kowalski.
\newblock \emph{Analytic Number Theory}.
\newblock American Mathematical Society Colloquium Publications, vol.~53,
  Providence, RI, 2004.

\bibitem{CojocMurty}
A.~C.~Cojocaru and M.~R.~Murty.
\newblock \emph{An Introduction to Sieve Methods and Their Applications}.
\newblock Cambridge University Press, 2005.

\bibitem{HalberstamRichert}
H.~Halberstam and H.-E.~Richert.
\newblock \emph{Sieve Methods}.
\newblock Academic Press, London, 1974.

\bibitem{Brun1919}
V.~Brun.
\newblock La s{\'e}rie $\frac{1}{5}+\frac{1}{7}+\cdots$ o\`u les d{\'e}nominateurs
  sont nombres premiers jumeaux est convergente ou finie.
\newblock \emph{Skrifter udgivne af Videnskabsselskabet i Christiania}, no.~3, 1919.

\bibitem{Maynard2015}
J.~Maynard.
\newblock Small gaps between primes.
\newblock \emph{Annals of Mathematics}, 181(1):383--413, 2015.

% BUG-P14-005 behoben: Label war Iwaniec1978, korrektes Jahr ist 1982 (JMSJ 34).
% Außerdem fehlt "the" im Journalnamen.
\bibitem{Iwaniec1982}
H.~Iwaniec.
\newblock On the Brun-Titchmarsh theorem.
\newblock \emph{Journal of the Mathematical Society of Japan}, 34:95--123, 1982.

\end{thebibliography}

\end{document}
